\chapter{کاوش گراف‌های شبکه‌های اجتماعی \texorpdfstring{\enparen{Social-Network Graphs}}{(Social-Network Graphs)}}
\label{ch:social-networks}

\section{مقدمه و ویژگی‌های توپولوژیک شبکه‌های اجتماعی}
با ظهور پلتفرم‌هایی نظیر لینکدین، توییتر و فیسبوک، گراف‌های شبکه‌های اجتماعی به یکی از مهم‌ترین منابع کلان‌داده تبدیل شده‌اند. در این گراف‌ها، گره‌ها بیانگر افراد (یا نهادها) و یال‌ها بیانگر روابط دوستی، تعامل، پیام‌رسانی یا همکاری هستند.

شبکه‌های اجتماعی در مقایسه با گراف‌های تصادفی ساده (مدل اردوش-رنی)، ویژگی‌های توپولوژیک منحصر‌به‌فردی دارند:
\begin{enumerate}
    \item \textbf{پدیده جهان کوچک \enparen{Small-World Phenomenon}}: میانگین طول کوتاه‌ترین مسیر میان هر دو فرد به طرز شگفت‌انگیزی کوتاه است (معروف به قانون «شش درجه جدایی»).
    \item \textbf{توزیع درجه قانون توانی}: توزیع درجات گره‌ها نامتقارن و دارای دنباله سنگین است؛ اکثر افراد روابط محدودی دارند، اما تعداد انگشت‌شماری «ابرگره» دارای میلیون‌ها ارتباط هستند.
    \item \textbf{ضریب خوشه‌بندی بالا \enparen{High Clustering Coefficient}}: اگر فرد $A$ با $B$ و $C$ دوست باشد، احتمال دوستی میان $B$ و $C$ بسیار بالاست (بسته‌شدن سه‌گانه\footnote{\lr{Triadic Closure}}).
    \item \textbf{ساختار جوامع \enparen{Community Structure}}: شبکه از گروه‌های متراکمی تشکیل شده است که درون هر گروه ارتباطات فراوان، اما میان گروه‌ها پیوندهای محدودی وجود دارد.
\end{enumerate}

\section{تشخیص جوامع و مرکزیت بینابینی یال‌ها}
چگونه می‌توان جوامع درون یک شبکه بزرگ را به صورت الگوریتمی کشف کرد؟
\begin{definition}[مرکزیت بینابینی یال \enparen{Edge Betweenness}]
مرکزیت بینابینی یک یال $e$ برابر است با تعداد کوتاه‌ترین مسیرها میان تمامی جفت‌گره‌های شبکه که از روی یال $e$ عبور می‌کنند:
\begin{equation}
C_B(e) = \sum_{s \ne t} \frac{\sigma_{st}(e)}{\sigma_{st}}
\end{equation}
که در آن $\sigma_{st}$ کل تعداد کوتاه‌ترین مسیرهای میان گره‌های $s$ و $t$، و $\sigma_{st}(e)$ تعداد مسیرهایی است که از یال $e$ می‌گذرند.
\end{definition}

\textbf{بینش شهودی}: یال‌هایی که درون یک جامعه قرار دارند دارای بینابینی پایینی هستند زیرا مسیرهای جایگزین متعددی درون جامعه وجود دارد؛ در مقابل، یال‌هایی که مانند یک «پل» دو جامعه مجزا را به هم متصل می‌کنند، دارای بالاترین مرکزیت بینابینی خواهند بود.

\subsection{الگوریتم گیروان-نیومن \texorpdfstring{\enparen{Girvan-Newman Algorithm}}{(Girvan-Newman Algorithm)}}
الگوریتم Girvan-Newman یک روش خوشه‌بندی سلسله‌مراتبی تقسیمی (بالا به پایین) است:

\begin{algorithmbox}[الگوریتم گیروان-نیومن \enparen{Girvan-Newman}]
\begin{enumerate}
    \item \textbf{محاسبه بینابینی تمام یال‌ها}: به ازای هر گره مبدأ $X$:
    \begin{itemize}
        \item \textbf{پیمایش رو به پایین \enparen{BFS}}: یک درخت جستجوی اول سطح \enparen{DAG} تشکیل دهید. به گره ریشه مقدار ۱ اختصاص دهید. به هر گره در سطوح بعدی، عددی برابر با مجموع مقادیر والدین مستقیم آن نسبت داده می‌شود (شمارش تعداد کوتاه‌ترین مسیرها از ریشه).
        \item \textbf{پیمایش رو به بالا (محاسبه جریان)}: از برگ‌های درخت آغاز کنید. هر گره یک واحد جریان به عنوان وزن خود در نظر می‌گیرد و جریان عبوری از هر یال ورودی به سمت والد $P$ متناسب با سهم والد در تعداد مسیرها $(\frac{\text{Score}(P)}{\text{Score}(Y)})$ محاسبه می‌شود. هر گره مجموع جریان‌های دریافتی از فرزندان را به علاوه ۱ کرده و به والدین بالادستی ارسال می‌نماید.
    \end{itemize}
    مجموع جریان‌های تمام درخت‌های BFS بر روی تمامی گره‌ها تقسیم بر ۲ (به دلیل بدون جهت بودن یال‌ها)، مقدار دقیق بینابینی هر یال را مشخص می‌کند.
    \item \textbf{حذف یال بحرانی}: یالی که دارای \textbf{بیشترین مرکزیت بینابینی} است از گراف حذف می‌گردد.
    \item \textbf{به‌روزرسانی و تکرار}: بینابینی یال‌های باقیمانده مجدداً محاسبه شده و فرآیند تکرار می‌شود تا ساختار خوشه‌ای مطلوب با بیشینه مدولاریتی حاصل شود.
\end{enumerate}
\end{algorithmbox}

\section{معیار پیمانه‌ای \texorpdfstring{\enparen{Modularity}}{(Modularity)}}
برای پاسخ به این پرسش که «تقسیم‌بندی شبکه به جوامع تا چه اندازه باکیفیت و معنادار است؟»، از شاخص \textbf{پیمانه‌ای}\footnote{\lr{Modularity}} $(Q)$ استفاده می‌شود.

\begin{definition}[پیمانه‌ای \enparen{Modularity}]
پیمانه‌ای کسر یال‌های درون جوامع را با کسر مورد انتظار همان یال‌ها در یک شبکه تصادفی با توزیع درجات یکسان مقایسه می‌کند:
\begin{equation}
Q = \frac{1}{2m} \sum_{i, j} \left( A_{ij} - \frac{k_i k_j}{2m} \right) \delta(c_i, c_j)
\end{equation}
که در آن:
\begin{itemize}
    \item $A_{ij}$: درایه ماتریس مجاورت (۱ اگر یال باشد، وگرنه ۰).
    \item $k_i, k_j$: درجه گره‌های $i$ و $j$.
    \item $m$: مجموع کل یال‌های گراف $(m = \frac{1}{2}\sum k_i)$.
    \item $c_i$: شناسه جامعه گره $i$.
    \item $\delta(c_i, c_j)$: تابع دلتا (۱ اگر هر دو گره در یک جامعه باشند، وگرنه ۰).
\end{itemize}
\end{definition}

مقدار $Q$ همواره در بازه $[-0.5, 1]$ قرار دارد. مقادیر مثبت بزرگ‌تر از $0.3$ نشان‌دهنده یک ساختار جامعه‌ای بسیار قوی و طبیعی در شبکه هستند.

\section{خوشه‌بندی طیفی گراف‌ها \texorpdfstring{\enparen{Spectral Clustering}}{(Spectral Clustering)}}
خوشه‌بندی طیفی یکی از قدرتمندترین ابزارهای جبر خطی در تحلیل گراف است که با استفاده از مقادیر و بردارهای ویژه ماتریس گراف عمل می‌کند.

\subsection{ماتریس لاپلاسین گراف \texorpdfstring{\enparen{Graph Laplacian}}{(Graph Laplacian)}}
برای گراف بدون جهت $G = (V, E)$:
\begin{enumerate}
    \item \textbf{ماتریس مجاورت $(A)$}: ماتریس متقارن با ابعاد $n \times n$.
    \item \textbf{ماتریس درجه $(D)$}: ماتریس قطری که درایه‌های قطر اصلی آن برابر با درجه هر گره است: $D_{ii} = d_i$.
    \item \textbf{ماتریس لاپلاسین گراف $(L)$}:
    \begin{equation}
    L = D - A
    \end{equation}
\end{enumerate}

\begin{theorem}[خواص بنیادین ماتریس لاپلاسین]
\begin{enumerate}
    \item ماتریس $L$ همواره متقارن و \textbf{نیمه‌معین مثبت}\footnote{\lr{Positive Semi-Definite}} است؛ یعنی تمامی مقادیر ویژه آن حقیقی و نامنفی هستند: $0 = \lambda_1 \le \lambda_2 \le \dots \le \lambda_n$.
    \item کوچک‌ترین مقدار ویژه همواره صفر است $(\lambda_1 = 0)$ و بردار ویژه متناظر با آن بردار تماماً یک است:
    \begin{equation}
    L \cdot \mathbf{1} = (D - A) \mathbf{1} = \mathbf{0}
    \end{equation}
    \item \textbf{تعداد مؤلفه‌های همبند}: چندگانگی\footnote{\lr{Multiplicity}} مقدار ویژه صفر دقیقاً برابر با تعداد مؤلفه‌های همبند گراف است.
\end{enumerate}
\end{theorem}

\subsection{بردار فیدلر \texorpdfstring{\enparen{Fiedler Vector}}{(Fiedler Vector)} و افراز گراف}
اگر گراف همبند باشد، دومین مقدار ویژه اکیداً مثبت است $(\lambda_2 > 0)$. بردار ویژه متناظر با $\lambda_2$، \textbf{بردار فیدلر}\footnote{\lr{Fiedler Vector}} نامیده می‌شود.

بردار فیدلر مسئله بهینه‌سازی پیوسته برش گراف را حل می‌کند:
\begin{equation}
\min_{x \perp \mathbf{1}, \|x\|=1} x^T L x = \min \sum_{(i, j) \in E} (x_i - x_j)^2
\end{equation}
اگر دو گره با یالی به هم متصل باشند، مؤلفه‌های متناظر آن‌ها در بردار فیدلر تمایل دارند مقادیری بسیار نزدیک به هم داشته باشند.

\begin{algorithmbox}[افراز گراف به کمک بردار فیدلر]
\begin{enumerate}
    \item ماتریس لاپلاسین $L = D - A$ را تشکیل دهید.
    \item دومین بردار ویژه کوچک‌تر $(\mathbf{x}_2)$ را محاسبه نمایید.
    \item گره‌ها را بر اساس مقدار درایه متناظرشان در $\mathbf{x}_2$ مرتب کنید.
    \item گره‌هایی را که درایه آن‌ها در $\mathbf{x}_2$ مثبت است در یک بخش $(A)$، و گره‌های با درایه منفی را در بخش دیگر $(B)$ قرار دهید.
\end{enumerate}
\end{algorithmbox}

\begin{summarybox}[جمع‌بندی فصل دهم]
\begin{itemize}
    \item شبکه‌های اجتماعی دارای ویژگی‌های ساختاری متمایز نظیر پدیده جهان کوچک، توزیع قانون توانی و خوشه‌بندی بالا هستند.
    \item الگوریتم گیروان-نیومن با حذف متوالی یال‌های با بالاترین بینابینی، ساختار جوامع را استخراج می‌کند.
    \item معیار پیمانه‌ای \enparen{Modularity} کیفیت افراز گراف را می‌سنجد.
    \item خوشه‌بندی طیفی با بهره‌گیری از ماتریس لاپلاسین $(L = D - A)$ و بردار فیدلر، برش بهینه گراف را در ابعاد طیفی میسر می‌سازد.
\end{itemize}
\end{summarybox}
